% Figura: mapa del capítulo sobre el pipeline.
% Se usa en la apertura del Capítulo 2. Cada componente lleva la sección que lo explica. Es la versión
% simplificada del diagrama del pipeline de sec:met_pipeline: si uno cambia, revisar el otro.
% Las secciones comunes a los tres componentes (fundamentos y atención) van en la leyenda y no en el
% diagrama: un recuadro que las agrupaba se superponía con la evaluación y no se entendía.

\begin{figure}[H]
    \centering
    \begin{tikzpicture}[x=1mm, y=1mm]

        \node[figentrada, text width=25mm] (img) at (0, 0) {Imagen\\dermatoscópica};

        \node[figtarea, text width=26mm] (clf) at (36, 14)
            {Clasificador\\{\footnotesize visión, §\ref{sec:marco_vision}}};
        \node[figtarea, text width=26mm] (knn) at (36, -14)
            {Codificador visual y kNN\\{\footnotesize §\ref{sec:marco_vlm}}};

        \node[figtarea, text width=26mm] (gen) at (88, 0)
            {Generador\\{\footnotesize lenguaje, §\ref{sec:marco_llm}\\VLM, §\ref{sec:marco_vlm}}};

        \node[figsalida, text width=14mm] (rep) at (117, 0) {Reporte};

        \node[fignodo, text width=29mm, draw=figpurpura!70!black, fill=figpurpura!12] (eval)
            at (113, -24) {Evaluación, §\ref{sec:marco_evaluacion}\\{\footnotesize generalización,
            §\ref{sec:marco_generalizacion}}};

        % Entrada: la imagen va a los dos componentes.
        \draw[figflecha] (img.east) -- ++(4, 0) |- (clf.west);
        \draw[figflecha] (img.east) -- ++(4, 0) |- (knn.west);

        % Salidas hacia el generador: tramo horizontal con la etiqueta encima y bajada hasta el generador.
        \draw[figflecha] (clf.east) -- node[figetiqueta, above, fill=none] {categoría} (67, 14)
            |- ([yshift=3mm]gen.west);
        \draw[figflecha] (knn.east) -- node[figetiqueta, below, fill=none, align=center] {casos\\parecidos} (67, -14)
            |- ([yshift=-3mm]gen.west);

        \draw[figflecha] (gen.east) -- (rep.west);
        \draw[figflecha] (rep.south) -- (rep.south |- eval.north);

    \end{tikzpicture}
    \caption[El capítulo recorrido sobre el \emph{pipeline}]{El capítulo recorrido sobre el
    \emph{pipeline}. Cada componente lleva la sección que presenta los conceptos que usa. El clasificador
    asigna una categoría a la imagen, un codificador visual recupera casos parecidos y el generador redacta
    el reporte con esa información; después la salida se evalúa y se examina si lo medido se sostiene
    fuera de los datos de entrenamiento. Los fundamentos del entrenamiento
    (Sección~\ref{sec:marco_fundamentos}) y la operación de atención (Sección~\ref{sec:marco_atencion}) son
    comunes a los tres componentes.}
    \label{fig:marco_mapa}
\end{figure}
